% Source: https://artofproblemsolving.com/community/c3h1385023_amc_10b_problem
% Stable identity: aops-topic-1385023-q23
Since $45=5\cdot9$ with $\gcd(5,9)=1$, the remainder $r$ of $N$ modulo $45$ is determined by its remainders modulo $5$ and modulo $9$ (Chinese remainder theorem).

\textbf{Modulo $5$:} $N$ ends in the digit $4$, so $N\equiv4\pmod5$.

\textbf{Modulo $9$:} a number is congruent to its digit sum modulo $9$, and writing $k$ in decimal contributes a digit sum congruent to $k$. So
\[
N\equiv1+2+\cdots+44=\dfrac{44\cdot45}{2}=990\equiv0\pmod9 .
\]
\textbf{Combine:} $r$ is a multiple of $9$ between $0$ and $44$, so $r\in\{0,9,18,27,36\}$, and $r\equiv4\pmod5$ leaves only $r=9$. The remainder is $\boxed{9}$, answer $\textbf{(C)}$.
